%%================================================
%% Filename: abstract.tex
%% Encoding: UTF-8
%% Author: Yuan Xiaoshuai - yxshuai@gmail.com
%% Created: 2012-04-24 00:21
%% Last modified: 2019-11-01 09:26
%%================================================
\begin{cabstract}

地方档案记录了街区名称、公共设施和居民生活的变化，是研究城市发展与社区治理的重要资料。档案数字化以后，研究者虽然可以通过关键词找到更多材料，却仍需逐份辨认同一地点在不同时期的名称，并核查一条记录究竟记述事实、转述旧闻，还是提出尚未实施的方案。本文以三个老城区的街道档案为例，研究怎样在保留原文和出处的前提下，把分散记录组织成可按地点与时间检索的资料，并评估这种组织方式对查找效率和证据判断的影响。

研究选取市政公报、街道工作记录和地方报纸三类材料，覆盖老城区公共服务设施调整较为频繁的二十年。资料整理以一条能在原文中定位的事件为单位，记录设施名称、行政区划、发生时间、事件类型和出处页码。遇到一页记载多项事件时分别拆分，同时为它们保留共同的出处；遇到多份材料记载同一事件时保留各自原文与日期。两名整理者先独立标注，再逐条讨论时间不清、地名重叠与设施分类不一致的材料，形成供后续检索对照的记录。

为处理地点名称的变化，本文建立带有有效时段的地名表，将历史名称和现在的名称连接到同一个地点标识，并记下支持这一对应关系的行政区划文件或地图。对于同名异地和改名年份不明的情形，系统保留可能对应的地点，提供原文线索供查询者判断。时间字段分开记录事件发生、信息发布和资料形成的日期，日期仅精确到月份时保留原有精度。这样可以把早年发生、后来才被报道的事项放回相应时段，并在需要时按报道日期检索。

在检索实验中，参与者使用同一批档案，分别完成关键词检索与基于地点和时间的检索任务。任务包括查找一处设施的历次搬迁、核对不同街区开始提供同类服务的时间，以及判断某则报道称的扩建是否已经完成。两种检索方式使用相同的档案影像与文字，参与者须在作答时附上出处。研究比较找到的相关记录数量、查找所需时间和出处的可核查程度，并记录参与者在查找失败时采用的补充线索。

结果显示，带有时段的地名表主要改善了涉及街道改名和边界调整的查找任务。使用同一地点标识后，旧名称与新名称下的记录能够一起列出，参与者更容易发现搬迁前后的联系。对设施名称长期稳定的任务，两种方式都能找到主要记录，增加地点筛选的收益较小。将事件日期和报道日期同时展示，有助于识别对旧事件的回顾。说明设施是否实际投入使用时，街道工作记录和后来公布的服务时间表往往需要结合核查。

对遗漏和误判的逐条复核表明，问题主要集中在地点描述过于简略的记录。例如，早期简报只写“车站附近”，相邻街区均有符合描述的地点，自动关联只能缩小查找范围。少数档案还用设施的计划名称指代后来更名的机构，必须结合前后文件确认其关系。为便于复核，检索界面同时显示匹配依据、原文位置和待确认的地点候选，允许使用者回到原页查看上下文。修订时保留原有记载，并另行登记对地点或日期的补充说明。后续使用者可以根据补充说明了解记录为何归到某条街道，也可以查看仍待确认的对应关系。对于材料本身没有提供足够线索的记录，系统保留未确定状态，并按档案形成单位与时间提供进一步浏览的入口。这让查询者能够分辨因资料缺失而无法作答的情况，以及尚可通过查阅其他材料解决的问题。

本文给出一套适用于地方档案的整理与检索方法，将地名沿革、事件时间和原文出处连在一起。对于需要跨年份核查公共设施变化的研究，这种方法能够减少重复查找，并使结论回到具体材料接受复核。它的应用效果取决于可利用的档案及地名沿革资料是否充分。将方法推广到其他地区时，应先整理当地常用地名的有效时段，核对区划变动，并对档案缺失较多的年份注明可检索的资料范围。

\end{cabstract}

\ckeywords{地方档案, 地名沿革, 时间检索, 公共服务设施, 证据核查}

\begin{eabstract}

Digital archives make records easier to search, but tracing a renamed facility still
requires examining the sources. Using three historic urban districts, this study
evaluates retrieval by place and date while preserving the original evidence.

We use municipal bulletins, neighborhood records, and newspapers spanning twenty
years. Two annotators identify public facility events and reconcile disagreements about
locations, dates, and facility types. Each report retains its own source and page.

A historical gazetteer connects old and new names and records when each was used.
Uncertain matches retain their source evidence. Event, publication, and document dates
are stored separately and retain their original precision.

Participants use keyword and place-based searches to trace relocations, compare when services
started, and verify reported expansions. We compare the relevant records retrieved,
search time, and whether answers can be checked against their cited sources.

The gazetteer mainly improves searches involving renaming or boundary changes.
Separate event and publication dates help identify retrospective accounts. Confirming
when a facility opened often requires checking both local records and service schedules.

Most errors arise from vague locations. The interface shows candidate locations, the
source page, and the evidence for each match. Corrections preserve the original
wording, and unresolved records can be browsed by source office and date.

Linking place-name histories, dates, and source pages reduces repeated searches and
helps verify public facility changes. Applying the method elsewhere requires checking
historical boundaries and identifying years with limited archival coverage.

\end{eabstract}

\ekeywords{Local archives, place-name history, temporal search,
public facilities, evidence checks}
